\documentclass[12pt,a4paper]{article}
\usepackage{ctex}
\usepackage{geometry}
\usepackage{enumitem}
\usepackage{hyperref}
\usepackage{titlesec}

% 页面边距
\geometry{left=2.5cm,right=2.5cm,top=2.5cm,bottom=2.5cm}

% 标题格式
\titleformat{\section}{\centering\bfseries\large}{\thesection}{1em}{}
\titleformat{\subsection}{\bfseries}{\thesubsection}{1em}{}

\begin{document}

\begin{center}
    \Large{\textbf{关于哈工大AI智能创新大赛使用\\“哈工大智能体开发平台”的通知}}
\end{center}

\vspace{1em}

各位参赛同学：

为更好地支撑本次AI智能创新大赛的开发工作，大赛组委会面向大家提供**“哈工大智能体开发平台”**作为本次比赛的开发与运行环境。

该平台基于火山HiAgent智能体开发框架建设，是学校统一部署的人工智能应用开发环境，集成了大模型调用、智能体编排、知识库管理及工具扩展等功能，可支持师生快速构建面向教学、科研和校园服务场景的AI智能体应用。通过平台提供的开发能力，参赛团队无需自行搭建复杂的AI开发环境，即可开展智能体设计、功能开发与应用测试。

现将平台开发权限申请方式说明如下：

\section{一、平台账号注册}
参赛师生访问智能体开发平台 \url{https://agent.hit.edu.cn}，使用个人统一身份认证账号完成登录。首次成功登录即视为完成注册，无需额外操作。

\section{二、申请开发权限}
完成平台注册后，请填写参赛人员开发权限申请：
\begin{center}
    \url{https://f.kdocs.cn/g/KFDlXN6B/}
\end{center}
平台管理员将在 24小时内 完成账号授权。

\section{三、确认授权结果}
授权完成后，再次登录智能体开发平台。如平台左侧菜单中出现智能体开发相关功能模块，则说明开发权限已成功开通。

\section{特别说明}
如您已具备校内智能体开发平台的开发权限，请务必仔细阅读以下文档，以规范、高效地开展后续开发工作：

\begin{itemize}
    \item 《智能体开发平台校内建设者标准操作流程参考（v2.0）》
\end{itemize}

该文档对平台使用规范、开发流程以及相关开发方法进行了详细说明，请各参赛团队在开发前认真查阅。

如在平台使用过程中遇到问题，可访问智能助手页面获取在线支持：
\begin{center}
    \url{https://agent.hit.edu.cn/product/llm/chat/d60l62ngvkeofbklliog}
\end{center}

后续大赛组委会将统一组织AI开发平台功能的使用培训，请随时关注群内通知。

\vspace{1.5em}

祝各位参赛选手在本次AI智能创新大赛中充分发挥创意与技术实力，取得优异成绩！

\end{document}